% FILE: 3-poster/parts/fonts.tex
% ==================================================================================================
% Typography (LuaLaTeX + fontspec). Two families realise the plan's hierarchy:
%   - Raleway          : display geometric sans -> title (Black) and headings (SemiBold/Bold)
%   - Source Sans Pro  : neutral sans          -> body, captions, labels, tables
% Fonts are resolved by name from the TeX Live tree (raleway, sourcesanspro packages).

\usepackage{fontspec}
\defaultfontfeatures{Ligatures=TeX}

% --- Body / default family: Source Sans Pro (sans) ---
% The ROMAN font is Latin Modern (serif) so that math alphabets (\mathrm, \mathbf) and the
% math "operators" font are serif, consistent with the serif math symbols. The document default
% is still sans (familydefault = sfdefault), so body text is Source Sans Pro.
\setmainfont{Latin Modern Roman}
\setsansfont{Source Sans Pro}
\renewcommand{\familydefault}{\sfdefault}

% --- Display family: Raleway (weights addressed explicitly) ---
% Title uses a lighter Medium weight (Black reads too heavy at poster scale).
\newfontfamily\titlefont{Raleway Medium}
\newfontfamily\headingfont{Raleway SemiBold}[
  BoldFont       = {Raleway Bold},
  ItalicFont     = {Raleway SemiBold Italic},
  BoldItalicFont = {Raleway Bold Italic}
]
\newfontfamily\headingbold{Raleway ExtraBold}

% --- Supporting body weights ---
\newfontfamily\bodylight{Source Sans Pro Light}
\newfontfamily\bodysemibold{Source Sans Pro Semibold}

% --- Mono family for section-heading kickers / metadata (per tests/claude-design-layout.tex) ---
\newfontfamily\kickerfont{Source Code Pro}

% --- Math text font: the body sans (Source Sans Pro), so descriptive words inside formulas
% (amplification, projection, in span, ...) match the body prose. Math symbols, the m/n vectors,
% matrices, and the "mix" indices keep their math font. Used via \mlab (see macros.tex). ---
\newfontfamily\mathtextfont{Source Sans Pro}

% --- Inline-math size correction. The body is Source Sans Pro but math symbols are Latin Modern
% (serif), whose smaller x-height makes inline states/matrices look smaller than the surrounding
% sans text. Enlarge the math size at the body text size (40 pt) so they visually match. (This only
% fires at 40 pt text, so the standalone supplement, which reuses this file at 11 pt, is unaffected.)
\DeclareMathSizes{40}{45}{32}{23}
% The same correction at the DISPLAY size used for the Approach's two headline formulae (the input and
% output vectors, block-approach move (a)). Without this line LaTeX has no math sizes for 54 pt text,
% falls back to defaults, and warns about missing font shapes. The brace readings under those formulae
% do NOT follow these sizes -- they are pinned by \brlab (macros.tex).
\DeclareMathSizes{54}{60}{40}{30}
